% ============================================================================
% C++ / Технологии программирования — общая преамбула конспекта
% Движок: XeLaTeX
% Основной стандарт примеров: C++17, если явно не указано иное.
% ============================================================================

% ---------- Язык, Unicode и шрифты ------------------------------------------
\usepackage{fontspec}
\usepackage{polyglossia}
\setdefaultlanguage{russian}
\setotherlanguage{english}

\defaultfontfeatures{Ligatures=TeX,Scale=MatchLowercase}
\setmainfont{DejaVu Serif}
\setsansfont{DejaVu Sans}
\setmonofont{DejaVu Sans Mono}

% ---------- Геометрия и базовая типографика --------------------------------
\usepackage[a4paper,top=20mm,bottom=22mm,left=22mm,right=18mm,headheight=15pt]{geometry}
\usepackage{microtype}
\usepackage{setspace}
\usepackage{parskip}
\usepackage{ragged2e}
\usepackage{enumitem}
\usepackage{titlesec}

\setlength{\parindent}{0pt}
\setlength{\parskip}{0.55em}
\setlist{nosep,leftmargin=2em}
\emergencystretch=2em

% ---------- Математика -------------------------------------------------------
\usepackage{amsmath}
\usepackage{amssymb}
\usepackage{mathtools}
\usepackage{amsthm}

% ---------- Цвета ------------------------------------------------------------
\usepackage{xcolor}

\definecolor{CppBlue}{HTML}{1F4E79}
\definecolor{CppDark}{HTML}{202124}
\definecolor{CppGray}{HTML}{5F6368}
\definecolor{CppLightGray}{HTML}{F4F5F7}
\definecolor{CppBorder}{HTML}{D0D7DE}
\definecolor{CppGreen}{HTML}{2E7D32}
\definecolor{CppOrange}{HTML}{B35C00}
\definecolor{CppRed}{HTML}{B42318}
\definecolor{CppPurple}{HTML}{6F42C1}
\definecolor{CppCyan}{HTML}{006D77}
\definecolor{CppYellowBg}{HTML}{FFF8DC}
\definecolor{CppRedBg}{HTML}{FFF1F0}
\definecolor{CppGreenBg}{HTML}{F1F8F3}
\definecolor{CppBlueBg}{HTML}{EEF6FF}
\definecolor{CppPurpleBg}{HTML}{F7F2FF}

% ---------- Рисунки, таблицы, плавающие объекты -----------------------------
\usepackage{graphicx}
\usepackage{eso-pic}
\usepackage{float}
\usepackage{booktabs}
\usepackage{tabularx}
\usepackage{longtable}
\usepackage{array}
\usepackage{multirow}
\usepackage{caption}
\usepackage{subcaption}

\graphicspath{{./}{figures/}{images/}}
\captionsetup{font=small,labelfont=bf}

\newcolumntype{Y}{>{\RaggedRight\arraybackslash}X}
\newcolumntype{C}[1]{>{\centering\arraybackslash}p{#1}}
\newcolumntype{L}[1]{>{\RaggedRight\arraybackslash}p{#1}}

% ---------- Ссылки -----------------------------------------------------------
\usepackage{hyperref}
\usepackage{xurl}
\usepackage[nameinlink,noabbrev]{cleveref}

\hypersetup{
  unicode=true,
  colorlinks=true,
  linkcolor=CppBlue,
  citecolor=CppGreen,
  urlcolor=CppPurple,
  pdfborder={0 0 0}
}

\urlstyle{same}

% ---------- Колонтитулы ------------------------------------------------------
\usepackage{fancyhdr}
\pagestyle{fancy}
\fancyhf{}
\fancyhead[L]{\small\sffamily Технологии программирования и C++}
\fancyhead[R]{\small\sffamily\nouppercase{\leftmark}}
\fancyfoot[C]{\small\sffamily\thepage}
\renewcommand{\headrulewidth}{0.4pt}
\renewcommand{\footrulewidth}{0pt}

% ---------- Код: listings ----------------------------------------------------
% listings выбран намеренно без shell-escape: PDF должен собираться обычным
% XeLaTeX. Основной код примеров — ASCII; русские пояснения даются вне листинга.
\usepackage{listings}
\usepackage{upquote}

\lstdefinestyle{codebase}{
  basicstyle=\ttfamily\small,
  backgroundcolor=\color{CppLightGray},
  frame=single,
  rulecolor=\color{CppBorder},
  framerule=0.5pt,
  framesep=7pt,
  xleftmargin=0.5em,
  xrightmargin=0.5em,
  breaklines=true,
  breakatwhitespace=false,
  postbreak=\mbox{\textcolor{CppGray}{$\hookrightarrow$}\space},
  columns=fullflexible,
  keepspaces=true,
  showstringspaces=false,
  showtabs=false,
  tabsize=4,
  captionpos=b,
  aboveskip=0.9em,
  belowskip=0.9em,
  escapeinside={(*@}{@*)},
  literate={~}{{\textasciitilde}}1
           {^}{{\textasciicircum}}1
}

\lstdefinestyle{cpp}{
  style=codebase,
  language=C++,
  keywordstyle=\bfseries\color{CppBlue},
  keywordstyle=[2]\bfseries\color{CppPurple},
  commentstyle=\itshape\color{CppGreen},
  stringstyle=\color{CppOrange},
  numberstyle=\scriptsize\color{CppGray},
  numbers=left,
  numbersep=9pt,
  morekeywords=[2]{override,final,noexcept,nullptr,constexpr,decltype,alignas,alignof,
                   static_assert,thread_local,concept,requires,co_await,co_yield,co_return}
}

\lstdefinestyle{cppplain}{
  style=cpp,
  numbers=none
}

\lstdefinestyle{shell}{
  style=codebase,
  language=bash,
  numbers=none,
  keywordstyle=\bfseries\color{CppBlue},
  commentstyle=\itshape\color{CppGreen},
  stringstyle=\color{CppOrange}
}

\lstdefinelanguage{CMake}{
  sensitive=false,
  morekeywords={cmake_minimum_required,project,add_executable,add_library,
    target_compile_features,target_compile_options,target_include_directories,
    target_link_libraries,set,option,if,elseif,else,endif,foreach,endforeach,
    function,endfunction,include,find_package,enable_testing,add_test,
    install,FetchContent_Declare,FetchContent_MakeAvailable},
  morecomment=[l]{\#},
  morestring=[b]"
}

\lstdefinestyle{cmake}{
  style=codebase,
  language=CMake,
  numbers=left,
  numberstyle=\scriptsize\color{CppGray},
  numbersep=9pt,
  keywordstyle=\bfseries\color{CppPurple},
  commentstyle=\itshape\color{CppGreen},
  stringstyle=\color{CppOrange}
}

\lstdefinelanguage{Diff}{
  morecomment=[f][\color{CppGreen}]{+},
  morecomment=[f][\color{CppRed}]{-},
  morecomment=[f][\color{CppCyan}]{@@}
}

\lstdefinestyle{diff}{
  style=codebase,
  language=Diff,
  numbers=none
}

\lstdefinestyle{output}{
  style=codebase,
  numbers=none,
  basicstyle=\ttfamily\footnotesize,
  backgroundcolor=\color{white},
  rulecolor=\color{CppBorder}
}

% Основные окружения для вставки кода.
\lstnewenvironment{cppcode}[1][]
  {\lstset{style=cpp,#1}}
  {}

\lstnewenvironment{cppcodeplain}[1][]
  {\lstset{style=cppplain,#1}}
  {}

\lstnewenvironment{shellcode}[1][]
  {\lstset{style=shell,#1}}
  {}

\lstnewenvironment{cmakecode}[1][]
  {\lstset{style=cmake,#1}}
  {}

\lstnewenvironment{outputcode}[1][]
  {\lstset{style=output,#1}}
  {}

\lstnewenvironment{diffcode}[1][]
  {\lstset{style=diff,#1}}
  {}

% Короткий inline-код. Использование: \cppinline|std::vector<int>|
\newcommand{\cppinline}{\lstinline[style=cppplain]}
\newcommand{\shellinline}{\lstinline[style=shell]}

% ---------- Цветные смысловые блоки -----------------------------------------
\usepackage[most]{tcolorbox}
\tcbuselibrary{breakable,skins,theorems,listings}

\newtcolorbox{importantbox}[1][Важно]{
  enhanced,breakable,
  colback=CppBlueBg,colframe=CppBlue,
  boxrule=0.7pt,arc=1.5mm,
  title={#1},fonttitle=\bfseries,
  left=1.2mm,right=1.2mm,top=1mm,bottom=1mm
}

\newtcolorbox{warningbox}[1][Осторожно]{
  enhanced,breakable,
  colback=CppYellowBg,colframe=CppOrange,
  boxrule=0.7pt,arc=1.5mm,
  title={#1},fonttitle=\bfseries,
  left=1.2mm,right=1.2mm,top=1mm,bottom=1mm
}

\newtcolorbox{ubbox}[1][Undefined behavior]{
  enhanced,breakable,
  colback=CppRedBg,colframe=CppRed,
  boxrule=0.8pt,arc=1.5mm,
  title={#1},fonttitle=\bfseries,
  left=1.2mm,right=1.2mm,top=1mm,bottom=1mm
}

\newtcolorbox{practicebox}[1][Практика]{
  enhanced,breakable,
  colback=CppGreenBg,colframe=CppGreen,
  boxrule=0.7pt,arc=1.5mm,
  title={#1},fonttitle=\bfseries,
  left=1.2mm,right=1.2mm,top=1mm,bottom=1mm
}

\newtcolorbox{performancebox}[1][Производительность]{
  enhanced,breakable,
  colback=CppPurpleBg,colframe=CppPurple,
  boxrule=0.7pt,arc=1.5mm,
  title={#1},fonttitle=\bfseries,
  left=1.2mm,right=1.2mm,top=1mm,bottom=1mm
}

\newtcolorbox{implementationbox}[1][Типичная реализация]{
  enhanced,breakable,
  colback=CppLightGray,colframe=CppGray,
  boxrule=0.7pt,arc=1.5mm,
  title={#1},fonttitle=\bfseries,
  left=1.2mm,right=1.2mm,top=1mm,bottom=1mm
}

\newtcolorbox{standardbox}[1][Что гарантирует стандарт]{
  enhanced,breakable,
  colback=white,colframe=CppCyan,
  boxrule=0.8pt,arc=1.5mm,
  title={#1},fonttitle=\bfseries,
  left=1.2mm,right=1.2mm,top=1mm,bottom=1mm
}

\newtcolorbox{whybox}[1][Почему?]{
  enhanced,breakable,
  colback=white,colframe=CppBlue,
  borderline west={2.2pt}{0pt}{CppBlue},
  boxrule=0pt,arc=0mm,
  title={#1},fonttitle=\bfseries,
  left=2mm,right=1mm,top=1mm,bottom=1mm
}

\newtcolorbox{mistakebox}[1][Типичная ошибка]{
  enhanced,breakable,
  colback=CppRedBg,colframe=CppRed,
  borderline west={2.2pt}{0pt}{CppRed},
  boxrule=0pt,arc=0mm,
  title={#1},fonttitle=\bfseries,
  left=2mm,right=1mm,top=1mm,bottom=1mm
}

% ---------- Определения, утверждения, леммы ---------------------------------
\theoremstyle{definition}
\newtheorem{definition}{Определение}[section]
\newtheorem{example}{Пример}[section]
\newtheorem{exercise}{Упражнение}[section]

\theoremstyle{plain}
\newtheorem{theorem}{Теорема}[section]
\newtheorem{lemma}{Лемма}[section]
\newtheorem{proposition}{Утверждение}[section]
\newtheorem{corollary}{Следствие}[section]

\theoremstyle{remark}
\newtheorem{remark}{Замечание}[section]

% ---------- Алгоритмы без дополнительной зависимости algorithm2e ------------
\newcounter{algorithmnote}[section]
\renewcommand{\thealgorithmnote}{\thesection.\arabic{algorithmnote}}

\newenvironment{algorithmnote}[1]{%
  \refstepcounter{algorithmnote}%
  \begin{tcolorbox}[
    enhanced,breakable,
    colback=white,colframe=CppDark,
    boxrule=0.7pt,arc=1.5mm,
    title={Алгоритм~\thealgorithmnote: #1},
    fonttitle=\bfseries,
    left=1.2mm,right=1.2mm,top=1mm,bottom=1mm]%
}{%
  \end{tcolorbox}%
}

% ---------- Небольшие семантические команды --------------------------------
\newcommand{\Cpp}{C\texttt{++}}
\newcommand{\CppStd}[1]{\textsf{C++#1}}
\newcommand{\term}[1]{\textbf{#1}}
\newcommand{\filepath}[1]{\texttt{#1}}
\newcommand{\flag}[1]{\texttt{#1}}
\newcommand{\command}[1]{\texttt{#1}}
\newcommand{\keyword}[1]{\texttt{#1}}
\newcommand{\type}[1]{\texttt{#1}}
\newcommand{\function}[1]{\texttt{#1}}
\newcommand{\headerfile}[1]{\texttt{<#1>}}

% Метки стандарта/реализации, чтобы визуально не смешивать уровни утверждений.
\newcommand{\standardtag}{\textsf{[standard]}}
\newcommand{\implementationtag}{\textsf{[implementation]}}
\newcommand{\practicetag}{\textsf{[practice]}}

% ---------- Оформление заголовков -------------------------------------------
\titleformat{\section}
  {\Large\bfseries\sffamily\color{CppDark}}
  {\thesection}{0.7em}{}

\titleformat{\subsection}
  {\large\bfseries\sffamily\color{CppBlue}}
  {\thesubsection}{0.7em}{}

\titleformat{\subsubsection}
  {\normalsize\bfseries\sffamily\color{CppDark}}
  {\thesubsubsection}{0.7em}{}

% ---------- Поведение содержания --------------------------------------------
\setcounter{tocdepth}{3}
\setcounter{secnumdepth}{3}

% ---------- cleveref: русские имена -----------------------------------------
\crefname{definition}{определение}{определения}
\Crefname{definition}{Определение}{Определения}
\crefname{theorem}{теорема}{теоремы}
\Crefname{theorem}{Теорема}{Теоремы}
\crefname{lemma}{лемма}{леммы}
\Crefname{lemma}{Лемма}{Леммы}
\crefname{proposition}{утверждение}{утверждения}
\Crefname{proposition}{Утверждение}{Утверждения}
\crefname{corollary}{следствие}{следствия}
\Crefname{corollary}{Следствие}{Следствия}
\crefname{example}{пример}{примеры}
\Crefname{example}{Пример}{Примеры}
\crefname{exercise}{упражнение}{упражнения}
\Crefname{exercise}{Упражнение}{Упражнения}
\crefname{algorithmnote}{алгоритм}{алгоритмы}
\Crefname{algorithmnote}{Алгоритм}{Алгоритмы}

% ============================================================================
% Конец общей преамбулы.
% ============================================================================


% ===== BEGIN ROBUST TOC NUMBER WIDTH FIX =====
% Ширины номеров задаются напрямую через стандартные LaTeX-макросы.
% Это не зависит от наличия пакета tocloft.
\makeatletter

% Верхний уровень: оставляем стандартный жирный стиль section,
% но увеличиваем место под номер секции, чтобы будущие "10", "11", ... не наезжали.
\renewcommand*\l@section[2]{%
  \ifnum \c@tocdepth >\z@
    \addpenalty\@secpenalty
    \addvspace{1.0em \@plus\p@}%
    \setlength\@tempdima{3.0em}%
    \begingroup
      \parindent \z@
      \rightskip \@pnumwidth
      \parfillskip -\@pnumwidth
      \leavevmode \bfseries
      \advance\leftskip\@tempdima
      \hskip -\leftskip
      #1\nobreak\hfil
      \nobreak\hb@xt@\@pnumwidth{\hss #2\kern-\p@}\par
    \endgroup
  \fi
}

% Подраздел: 3.6em вместо стандартного узкого поля.
% Этого достаточно для номеров 1.10, 12.18 и подобных.
\renewcommand*\l@subsection{%
  \@dottedtocline{2}{1.5em}{3.6em}%
}

% Подподраздел: номер может быть вида 1.3.10, поэтому поле ещё шире.
% Отступ начинается после ширины subsection, чтобы уровни визуально не слипались.
\renewcommand*\l@subsubsection{%
  \@dottedtocline{3}{5.1em}{4.6em}%
}

\makeatother
% ===== END ROBUST TOC NUMBER WIDTH FIX =====
